\section{Conclusion}

This thesis studied deception detection in multilingual (English/Bangla/Banglish) computer-mediated communication under a privacy-preserving federated learning regime, using a shared attention-pooling classification architecture built on top of five interchangeable multilingual and Bangla-aware pretrained Transformer encoders. Unlike most prior federated fake-news and misinformation detection work, which partitions clients by user or by random subsample of a single task, the five federated clients in this thesis each held an entire, non-overlapping deception domain -- fake news, product reviews, phishing, political statements, and job scams -- producing a deliberately severe form of non-\gls{fl}-friendly statistical heterogeneity, quantitatively confirmed (\Cref{sec:dataset-heterogeneity}) to be over twenty times more divergent in client vocabulary than a random partition of the same data, while label and language distributions remain balanced across clients. This severity of heterogeneity motivated the use of the FedProx optimization strategy in place of plain FedAvg.

The final results support the following headline conclusions:
\begin{itemize}
    \item \textbf{Best backbone overall:} BanglishBERT, under both training paradigms -- 82.9\% accuracy centrally, 76.0\% accuracy federatedly, ahead of MuRIL, DistilXLM-R, mBERT, and ALBERT Multilingual on every metric evaluated.
    \item \textbf{Accuracy-retention, all five backbones:} every backbone now has a matched centralized/federated pair, and every one retains 89.7--94.3\% of its centralized accuracy and macro-F1 under federation (\Cref{sec:retention}) -- a narrow, consistent band with no anomalous backbone.
    \item \textbf{Source of the pooled-accuracy gap:} the per-domain breakdown (\Cref{sec:domain-breakdown}) shows this cost is not uniform: Fake News and Phishing are handled well by every backbone (85--90\% average accuracy, competitive with single-domain prior work), while Product Reviews and Political Statements are weak across every backbone regardless of architecture, indicating a domain-level difficulty rather than a modeling weakness.
    \item \textbf{Overall verdict on the privacy-utility trade-off:} federated training is a practical, low-cost alternative to centralized training for this task. Given the severity of the client heterogeneity introduced by the five-domain partitioning (\Cref{sec:dataset-heterogeneity}), a consistent 6--10\% accuracy cost across all five backbones is a favorable outcome, not a weak one.
\end{itemize}

Taken together, these results support the conclusion that combining a Bangla/Indian-language-specialized encoder (BanglishBERT or MuRIL) with FedProx-based federated training is a practical way to build a multilingual deception detector without centralizing sensitive user text, retaining the large majority of centralized accuracy despite an unusually adversarial, task-partitioned, and quantifiably more heterogeneous client setting than has been evaluated in prior federated fake-news detection work.

\section{Future Work}
\label{sec:future-work}

Several directions remain open for extending this work:

\begin{enumerate}
    \item \textbf{FedProx hyperparameter tuning.} Systematically evaluate the effect of the proximal weight $\mu$ (beyond the default 0.01, e.g.\ in the 0.1--1.0 range) and the balance between local epochs per round and total communication rounds (e.g.\ \texttt{LOCAL\_EPOCHS}~=~1 with \texttt{FED\_ROUNDS}~=~20), which \Cref{sec:federated} identifies as a plausible way to further reduce client drift under this thesis's especially severe task-partitioned client setting.
    \item \textbf{Root-cause analysis of weak domains.} Product Reviews and Political Statements are weak across every backbone tested (\Cref{sec:domain-breakdown}); further work is needed to determine how much of this is intrinsic task difficulty (e.g.\ Political Statements requiring external world knowledge, as in the LIAR benchmark \cite{wang2017liar}) versus an artifact of the translation/transliteration process used to construct the Bangla and Banglish portions of the dataset (\Cref{chap:dataset}).
    \item \textbf{Alternative federated aggregation strategies.} Compare FedProx against other approaches designed for heterogeneous clients, such as personalized federated learning or clustered federated learning \cite{ghosh2020efficient}, which may be particularly well suited to a setting where clients are heterogeneous by design (each holding a distinct task) rather than by incidental sampling noise -- especially given the 21.85$\times$ vocabulary divergence quantified in \Cref{sec:dataset-heterogeneity}.
    \item \textbf{Explainability.} Extend the attention-pooling weights already produced by the architecture in \Cref{sec:arch} into a systematic explainability analysis, in the spirit of Ilias et al.\ \cite{ilias2022explainable}, to verify that each backbone is relying on genuinely deception-relevant linguistic signal rather than domain- or language-specific artifacts introduced by the translation/transliteration process described in \Cref{chap:dataset}.
    \item \textbf{Beyond text.} Incorporate additional modalities where available (e.g.\ metadata, network/propagation signal for fake news, or sender reputation for phishing), following the general direction of multimodal deception-detection work such as Gallardo-Antol{\'i}n and Montero \cite{gallardo2021detecting}, while preserving the federated, privacy-preserving training regime.
    \item \textbf{Real-world, non-synthetic multilingual data.} Replace the translated/transliterated Bangla and Banglish portions of the dataset (\Cref{chap:dataset}) with naturally-occurring Bangla and Banglish deceptive text collected directly from Bangladeshi digital platforms, to validate that the results reported in \Cref{chap:results} are not an artifact of machine-translation style.
\end{enumerate}
